\documentclass[11pt]{article}

\usepackage[margin=1in]{geometry}
\usepackage[T1]{fontenc}
\usepackage[utf8]{inputenc}
\usepackage{lmodern}
\usepackage{microtype}
\usepackage{amsmath,amssymb,amsthm,mathtools}
\usepackage{bm}
\usepackage{booktabs}
\usepackage{array}
\usepackage{enumitem}
\usepackage[numbers,sort&compress]{natbib}
\usepackage[colorlinks=true,linkcolor=blue!60!black,citecolor=blue!60!black,urlcolor=blue!60!black]{hyperref}

\newtheorem{definition}{Definition}[section]
\newtheorem{proposition}[definition]{Proposition}
\newtheorem{lemma}[definition]{Lemma}
\newtheorem{corollary}[definition]{Corollary}
\theoremstyle{remark}
\newtheorem{remark}[definition]{Remark}
\newtheorem{example}[definition]{Example}

\newcommand{\R}{\mathbb{R}}
\newcommand{\C}{\mathbb{C}}
\newcommand{\HH}{\mathbb{H}}
\newcommand{\SU}{\mathrm{SU}}
\newcommand{\U}{\mathrm{U}}
\newcommand{\Id}{I}
\newcommand{\ii}{\mathbf{i}}
\newcommand{\jj}{\mathbf{j}}
\newcommand{\kk}{\mathbf{k}}
\newcommand{\norm}[1]{\left\lVert #1 \right\rVert}
\newcommand{\set}[1]{\left\{ #1 \right\}}
\newcommand{\uoneq}{\mathrm{u1q}}
\newcommand{\normalize}{\operatorname{Norm}}
\newcommand{\sgncanon}{\operatorname{SignCan}}
\newcommand{\canon}{\operatorname{Can}}
\newcommand{\Fuse}{\operatorname{Fuse}}
\newcommand{\Recon}{\operatorname{Recon}}
\newcommand{\Baseline}{\operatorname{Base}}
\newcommand{\GateCount}{\operatorname{G}}
\newcommand{\Depth}{\operatorname{D}}
\newcommand{\Eq}{\operatorname{Eq}}

\title{\textbf{Quaternionic Gate Fusion on $S^3$:}\\
Geometry-Native Optimization for Single-Qubit Circuit Segments}

\author{RQM Technologies Research Team\thanks{Organizational author record used for this draft. Correspondence: \texttt{research@rqm-technologies.com}.}\\
RQM Technologies}

\date{Version 0.1.0 \\ Draft dated 2026-04-09}

\begin{document}
\maketitle

\begin{abstract}
This paper studies a local optimization pass for single-qubit circuit segments
using the canonical quaternionic intermediate representation defined in
\texttt{qqc-002-canonical-ir}. A maximal run of adjacent single-qubit gates on
one wire is translated into unit quaternions, multiplied in circuit order,
re-normalized, and sign-canonicalized on the quotient $q\sim -q$ induced by
global-phase removal. The result is a unique canonical fused representative in
$S^3$ for the segment action modulo global phase. A layered reconstruction
strategy then emits either no gate, an exact named gate, an exact axis-aligned
rotation, or a generic single-qubit primitive when no simpler interpretable
output is available. The paper proves exact equivalence preservation up to
global phase, determinism of the fused representative under the fixed
canonicalization rule, and minimality of the fully fused segment in the
quaternionic IR itself. A small reference benchmark corpus shipped with the
paper package is evaluated by an included plain-Python harness. On 13 measured
segments containing 36 input single-qubit gates, the local fusion pass
reconstructs to 11 output gates while preserving exact segment action under the
authoritative canonical-quaternion equivalence check. Two segments eliminate to
identity, and 8 of 13 outputs are recovered as either named gates or
axis-aligned rotations. These measurements are limited to the included
single-qubit reference corpus and are not presented as a claim of universal
compiler superiority.
\end{abstract}

\noindent\textbf{Keywords:} quaternionic gate fusion; single-qubit optimization; compiler method; $\SU(2)$; unit quaternions; local circuit simplification

\section{Introduction}

Paper 1 in this series, \texttt{qqc-001-foundations}, established the fixed
quaternion--$\SU(2)$ convention used throughout the repository. Paper 2,
\texttt{qqc-002-canonical-ir}, defined a canonical single-qubit intermediate
representation $\uoneq(w,x,y,z)$ together with deterministic normalization and
sign-canonicalization rules. The present paper asks the next compiler-facing
question: once the representation exists, what optimization rule does it induce?

The answer is local and operational. If a circuit contains a maximal run of
adjacent single-qubit gates on one wire, then those gates may be fused by
quaternion multiplication. This is not an analogy or a heuristic. It is the
group law of the representation chosen in papers 1 and 2. The optimizer maps
the segment into quaternions, multiplies in circuit order, re-canonicalizes the
product, and then reconstructs an interpretable output whenever possible.

This paper is therefore not a manifesto about quaternions. It is a method paper
about a local optimization pass. The pass is limited to single-qubit segments
between barriers such as entangling gates, measurement, reset, or classical
control boundaries. It does not optimize entangling gates, and it does not
claim to solve full-circuit synthesis or compiler optimality.

Three contributions are the focus:
\begin{enumerate}[label=(\arabic*), leftmargin=1.5em]
\item a formally specified fusion operator on single-qubit segments;
\item a reconstruction layer that is exact when identity, named-gate, or
axis-aligned forms are available, but otherwise falls back honestly to a
generic one-qubit primitive; and
\item a measured evaluation on a shipped reference corpus with a clearly stated
syntax-preserving baseline.
\end{enumerate}

\section{Background and dependency on papers 1--2}

\subsection{The inherited mathematical model}

We adopt without change the convention fixed in paper 1:
\begin{equation}
\Phi(a+b\ii+c\jj+d\kk)
=
\begin{pmatrix}
a-di & -c-bi \\
c-bi & a+di
\end{pmatrix}
=
a\Id-i(b\sigma_x+c\sigma_y+d\sigma_z).
\label{eq:phi}
\end{equation}
Under this convention, unit quaternions model $\SU(2)$ exactly, and single-qubit
gates in $\U(2)$ admit representatives in $\SU(2)$ after removal of global
phase.

\begin{proposition}[Single-qubit gate actions admit unit-quaternion representatives]
For every single-qubit gate $U\in\U(2)$ there exists a unit quaternion
$q\in\HH$ such that $\Phi(q)$ represents the same single-qubit action as $U$ on
pure states modulo global phase. The representative $q$ is unique up to sign.
\end{proposition}

\begin{proof}
By paper 1, every $U\in\U(2)$ admits a representative $V\in\SU(2)$ after
removing determinant phase, and $V$ is unique up to sign. The restriction of
$\Phi$ to the unit quaternions is an isomorphism onto $\SU(2)$, so $V$
corresponds to a unique sign pair $\set{q,-q}$.
\end{proof}

\subsection{The inherited canonical IR}

Paper 2 defined the canonical single-qubit IR object
\[
\uoneq(w,x,y,z),
\]
interpreted as the unit quaternion
\[
q=w+x\ii+y\jj+z\kk,\qquad w^2+x^2+y^2+z^2=1.
\]
It also fixed the canonicalization map
\[
\canon=\sgncanon\circ\normalize
\]
with primary rule $w>0$ and an equatorial tie-break when $w=0$.

\begin{remark}
The present paper depends conceptually on paper 2 for determinism. Quaternion
multiplication alone would not choose a unique stored representative, because
$q$ and $-q$ encode the same single-qubit action modulo global phase. Fusion
therefore uses multiplication \emph{followed by} canonicalization, not raw
multiplication alone.
\end{remark}

\section{Problem setup}

\begin{definition}[Single-qubit segment]
A \emph{single-qubit segment} on wire $a$ is a maximal consecutive run of gates
that act only on wire $a$, bounded by the beginning or end of the circuit, or
by a boundary operation that lies outside the present single-qubit model for
that wire. Typical boundaries include entangling gates, measurement, reset, and
classically controlled branching points.
\end{definition}

\begin{definition}[Target optimization problem]
Given a single-qubit segment
\[
\mathrm{seg}=(g_1,g_2,\dots,g_n),
\]
find a replacement segment $\mathrm{seg}'$ such that:
\begin{enumerate}[label=(\alph*), leftmargin=1.5em]
\item $\mathrm{seg}'$ represents exactly the same single-qubit action as
$\mathrm{seg}$ modulo global phase;
\item $\mathrm{seg}'$ is deterministic under the fixed canonicalization
convention; and
\item $\GateCount(\mathrm{seg}')$ and $\Depth(\mathrm{seg}')$ are no larger than
their baseline values for the original segment.
\end{enumerate}
\end{definition}

\begin{remark}
For a purely sequential single-qubit segment, gate count and segment-local depth
coincide. In a larger circuit, however, this paper claims reduction only for the
single-qubit run itself. Entangling depth is not changed by the present method.
\end{remark}

\begin{definition}[Authoritative equivalence criterion]
For two single-qubit segments $\mathrm{seg}$ and $\mathrm{seg}'$, define
\[
\Eq(\mathrm{seg},\mathrm{seg}')
\]
to hold when their canonical fused quaternions agree exactly:
\[
\canon(q_n\cdots q_1)=\canon(q'_m\cdots q'_1),
\]
where the $q_i$ and $q'_j$ are the translated unit-quaternion representatives of
the gates in application order.
\end{definition}

\section{Quaternionic fusion method}

\subsection{Translation and multiplication order}

Let
\[
\mathrm{seg}=(g_1,g_2,\dots,g_n)
\]
be a single-qubit segment written in circuit application order. For each gate
$g_k$, let $q_k$ be its canonical unit-quaternion representative under the map
of paper 2.

\begin{definition}[Quaternionic fusion operator]
Define the fused quaternionic payload of the segment by
\[
\Fuse(\mathrm{seg})=\canon(q_n q_{n-1}\cdots q_1).
\label{eq:fuse}
\]
\end{definition}

\begin{remark}[Order convention]
Equation \eqref{eq:fuse} is intentionally right-to-left in the gate index. The
first gate applied to the state is $g_1$, and under \eqref{eq:phi} one has
\[
\Phi(q_n q_{n-1}\cdots q_1)
=
\Phi(q_n)\Phi(q_{n-1})\cdots\Phi(q_1),
\]
which is the matrix product for sequential application in that order.
\end{remark}

\subsection{Normalization and sign-canonicalization}

The fusion pass uses the canonicalization rule inherited from paper 2:
\begin{enumerate}[label=(\arabic*), leftmargin=1.5em]
\item translate each gate into a canonical unit quaternion;
\item multiply in circuit order as in \eqref{eq:fuse};
\item re-normalize if numerical drift has occurred; and
\item apply sign-canonicalization with the primary rule $w>0$ and the
equatorial tie-break when $w=0$.
\end{enumerate}

\begin{definition}[Baseline flow]
The benchmark baseline used in this paper is the
\emph{syntax-preserving baseline}
\[
\Baseline(\mathrm{seg})=\mathrm{seg},
\]
meaning that no fusion is performed and the gate list is left unchanged.
\end{definition}

\section{Axis-aware reconstruction}

The fused quaternion is a complete canonical segment representative, but a
compiler often benefits from an interpretable output layer.

\begin{definition}[Axis-aware reconstruction]
Given a canonical fused quaternion $q$, the reconstruction operator
\[
\Recon(q)
\]
applies the following layered rule:
\begin{enumerate}[label=(\arabic*), leftmargin=1.5em]
\item if $q=(1,0,0,0)$, emit no gate;
\item if $q$ matches a stored exact canonical tuple for a named gate such as
$X$, $Y$, $Z$, $H$, $S$, $S^\dagger$, $T$, or $T^\dagger$, emit that named gate;
\item if exactly one of $(x,y,z)$ is nonzero, recover the corresponding exact
axis-aligned rotation $R_x(\theta)$, $R_y(\theta)$, or $R_z(\theta)$;
\item otherwise emit a generic canonical one-qubit primitive
$\uoneq(w,x,y,z)$ or an equivalent internal opaque primitive.
\end{enumerate}
\end{definition}

\begin{remark}
The reconstruction layer is not the equivalence proof. Equivalence is certified
by canonical quaternion equality. Reconstruction is a presentation and compiler
interface layer on top of that exact representation.
\end{remark}

\begin{proposition}[Exact axis-aligned cases admit direct recovery]
If the canonical fused quaternion has exactly one nonzero imaginary component,
then it reconstructs exactly to an axis-aligned rotation family member
$R_x(\theta)$, $R_y(\theta)$, or $R_z(\theta)$.
\end{proposition}

\begin{proof}
If only one of $(x,y,z)$ is nonzero, then the quaternion already has
axis--angle form about a coordinate axis:
\[
q=\cos(\theta/2)+\sin(\theta/2)\ii,
\quad
q=\cos(\theta/2)+\sin(\theta/2)\jj,
\quad \text{or} \quad
q=\cos(\theta/2)+\sin(\theta/2)\kk.
\]
Under \eqref{eq:phi}, these are exactly $R_x(\theta)$, $R_y(\theta)$, and
$R_z(\theta)$.
\end{proof}

\begin{remark}[Named-gate recovery is not globally unique]
The canonical fused quaternion is unique under the fixed convention. A
human-readable named-gate decomposition generally is not. The present paper only
claims uniqueness for the quaternionic representative and exact recovery for the
special cases explicitly recognized by $\Recon$.
\end{remark}

\section{Theoretical properties}

\begin{proposition}[Fusion preserves exact segment action up to global phase]
For every single-qubit segment $\mathrm{seg}$,
\[
\Eq(\mathrm{seg},\Fuse(\mathrm{seg}))
\]
holds when $\Fuse(\mathrm{seg})$ is interpreted as its fused canonical
representative.
\end{proposition}

\begin{proof}
Let $q_1,\dots,q_n$ be the gate representatives in application order. Under
\eqref{eq:phi}, the matrix product of the segment is
\[
\Phi(q_n)\cdots\Phi(q_1)=\Phi(q_n\cdots q_1).
\]
Canonicalization does not change the represented segment action modulo global
phase, because it only normalizes and chooses the unique representative from the
sign pair $\set{q,-q}$. Hence the fused canonical quaternion represents exactly
the same single-qubit action as the original segment.
\end{proof}

\begin{proposition}[Determinism of the fused representative]
For every nonempty single-qubit segment $\mathrm{seg}$, the value
$\Fuse(\mathrm{seg})$ is uniquely determined by the segment action and the fixed
canonicalization convention of paper 2.
\end{proposition}

\begin{proof}
The translated product $q_n\cdots q_1$ is uniquely determined by the segment
action up to the sign ambiguity inherited from $\SU(2)$. The map $\canon$
chooses exactly one representative from each sign class. Therefore the fused
representative is unique.
\end{proof}

\begin{corollary}[Segment replacement by one canonical payload]
Every nonidentity single-qubit segment may be replaced by one canonical
quaternionic payload without changing the represented segment action modulo
global phase. Identity segments may be replaced by the empty segment.
\end{corollary}

\begin{proof}
This follows immediately from the previous two propositions.
\end{proof}

\begin{proposition}[Minimality in the quaternionic IR]
Within the quaternionic IR itself, the fully fused representation of a nonidentity
single-qubit segment has minimal gate count equal to one.
\end{proposition}

\begin{proof}
The quaternionic IR stores segment action as a single canonical payload
$\uoneq(w,x,y,z)$. Any nonidentity segment action has at least one such payload,
and no smaller nonempty representation exists inside that IR. Identity is the
unique case in which the representation can be reduced to zero emitted gates.
\end{proof}

\section{Benchmark methodology}

\subsection{Reference corpus}

The benchmark corpus shipped with this paper package is deliberately modest. It
is intended as a reference evaluation of the local fusion method rather than a
claim of production-scale compiler performance. The included harness evaluates
13 segments across five categories:
\begin{enumerate}[label=(\arabic*), leftmargin=1.5em]
\item textbook single-qubit chains;
\item rotation-heavy segments;
\item mixed named-gate segments;
\item single-qubit runs extracted from small two-qubit circuits by entangling
boundaries; and
\item an ansatz-style single-wire segment.
\end{enumerate}

\subsection{Metrics}

For each segment $\mathrm{seg}$, the harness records:
\begin{align*}
\GateCount_{\mathrm{orig}}(\mathrm{seg}) &= \text{original number of gates},\\
\GateCount_{\mathrm{opt}}(\mathrm{seg}) &= \text{number of reconstructed output gates},\\
\Depth_{\mathrm{orig}}(\mathrm{seg}) &= \text{original segment-local depth},\\
\Depth_{\mathrm{opt}}(\mathrm{seg}) &= \text{optimized segment-local depth}.
\end{align*}
It also records:
\begin{itemize}[leftmargin=1.5em]
\item whether the reconstructed output is identity, named, axis-aligned, or
generic;
\item whether the reconstructed output passes the authoritative equivalence check
after re-translation into canonical quaternion form; and
\item whether the segment collapsed to one canonical representative before
reconstruction.
\end{itemize}

\subsection{Baseline}

The baseline is the syntax-preserving flow $\Baseline$. It performs no fusion and
leaves the segment unchanged. The comparison is therefore intentionally modest:
the paper measures what the local quaternionic pass does relative to doing
nothing on those segments.

\begin{remark}
No comparison is claimed here against a major external compiler stack. The
present goal is to study a local method and to ship a reproducible reference
harness for it.
\end{remark}

\section{Worked examples and measured reference-corpus results}

\subsection{Worked examples}

\begin{example}[$R_x(\pi/2);R_x(\pi/2)$]
The translated quaternions are
\[
q_1=q_2=\left(\frac{\sqrt2}{2},\frac{\sqrt2}{2},0,0\right).
\]
Hence
\[
\Fuse(\mathrm{seg})=\canon(q_2 q_1)=(0,1,0,0).
\]
Axis-aware recovery first identifies the $x$ axis, yielding $R_x(\pi)$, and the
named-gate table can further recognize that canonical tuple as $X$. The segment
changes from gate count and local depth $2$ to $1$.
\end{example}

\begin{example}[$R_z(\pi/4);R_z(\pi/4)$]
The fused quaternion is
\[
\canon\!\left(
\left(\cos\frac{\pi}{8},0,0,\sin\frac{\pi}{8}\right)^2
\right)
=
\left(\frac{1}{\sqrt2},0,0,\frac{1}{\sqrt2}\right).
\]
This is axis-aligned about $z$ and is also recognized as the named gate $S$.
Again the segment changes from $(2,2)$ in count and local depth to $(1,1)$.
\end{example}

\begin{example}[$H;Z;H$]
Using the canonical representatives of paper 2,
\[
q_H=\left(0,\frac{1}{\sqrt2},0,\frac{1}{\sqrt2}\right),
\qquad
q_Z=(0,0,0,1).
\]
Fusion gives
\[
\Fuse(\mathrm{seg})=\canon(q_H q_Z q_H)=(0,1,0,0),
\]
which reconstructs to $X$. This is the expected similarity transformation, now
expressed directly as a quaternionic local optimization.
\end{example}

\begin{example}[A segment between two entangling boundaries]
Consider the wire-0 run
\[
T^\dagger;R_x(\pi/4)
\]
extracted from a larger two-qubit circuit between two CNOT barriers. Fusion
gives
\[
\Fuse(\mathrm{seg})\approx
(0.8535534,\,0.3535534,\,0.1464466,\,-0.3535534).
\]
This output is not axis-aligned and does not match the exact named-gate lookup
table, so reconstruction emits a generic $\uoneq$ primitive. Even in this case,
the local optimization still reduces the segment from two gates to one while
preserving exact segment action modulo global phase.
\end{example}

\subsection{Measured reference-corpus results}

The included benchmark harness in
\texttt{artifacts/code/benchmark\_gate\_fusion.py} computes the measured results
stored in \texttt{artifacts/code/benchmark\_results.json}. On the shipped
reference corpus:
\begin{itemize}[leftmargin=1.5em]
\item 13 segments were evaluated;
\item original gate count was $36$ and optimized gate count was $11$;
\item original segment-local depth was $36$ and optimized segment-local depth
was $11$;
\item 2 segments collapsed to identity and emitted no gate;
\item all 13 segments collapsed to a single canonical quaternionic
representative before reconstruction; and
\item all 13 equivalence checks passed.
\end{itemize}

These numbers correspond to a gate-count and segment-depth reduction of
\[
36-11=25,
\]
or approximately $69.4\%$ relative to the syntax-preserving baseline on this
reference corpus.

The measured reconstruction breakdown is:
\begin{center}
\begin{tabular}{@{}lr@{}}
\toprule
Reconstruction kind & Count \\
\midrule
named gate (including identity) & 6 \\
axis-aligned rotation & 2 \\
generic one-qubit primitive & 5 \\
\bottomrule
\end{tabular}
\end{center}
Thus 8 of 13 outputs (about $61.5\%$) are recoverable as named or exact
axis-aligned forms, while the remaining 5 are emitted as generic canonical
primitives.

Category-level results are shown in Table~\ref{tab:category-results}.

\begin{table}[h]
\centering
\begin{tabular}{@{}lrrrr@{}}
\toprule
Category & Segments & Original gates & Optimized gates & Reduction \\
\midrule
Textbook single-qubit chains & 3 & 7 & 3 & 4 \\
Mixed named-gate segments & 3 & 9 & 2 & 7 \\
Rotation-heavy segments & 4 & 11 & 3 & 8 \\
Segments from two-qubit circuits & 2 & 4 & 2 & 2 \\
Ansatz-style segments & 1 & 5 & 1 & 4 \\
\midrule
Total & 13 & 36 & 11 & 25 \\
\bottomrule
\end{tabular}
\caption{Measured gate-count results on the shipped reference corpus under the
syntax-preserving baseline. Because every benchmark item is a sequential
single-qubit segment, the same totals also describe segment-local depth.}
\label{tab:category-results}
\end{table}

\begin{remark}
These measurements are intentionally narrow. They demonstrate that the local
fusion rule is operational and measurable on a shipped corpus. They do not
establish universal superiority over other compiler flows, and they do not
address whole-circuit optimality.
\end{remark}

\section{Standard mathematics vs.\ RQM Technologies framing}

\subsection{Standard ingredients}

The following ingredients are standard:
\begin{enumerate}[label=(\arabic*), leftmargin=1.5em]
\item single-qubit gates live in $\U(2)$ and admit $\SU(2)$ representatives
modulo global phase;
\item unit quaternions model $\SU(2)$ under the fixed convention \eqref{eq:phi};
\item sequential single-qubit gate composition corresponds to quaternion
multiplication; and
\item $q\sim -q$ is the residual ambiguity after removal of global phase.
\end{enumerate}

\subsection{Project-specific contribution}

The project-specific contribution of this paper is the optimization method:
\begin{enumerate}[label=(\alph*), leftmargin=1.5em]
\item use the canonical $\uoneq$ representation of paper 2 as the working IR;
\item define a local fusion pass by quaternion multiplication plus
re-canonicalization;
\item use a layered reconstruction rule for identity, named-gate, axis-aligned,
and generic outputs; and
\item evaluate the pass against a syntax-preserving baseline on a shipped
reference corpus.
\end{enumerate}

These are compiler-method choices. They are not new mathematics, and they are
not claims about hardware advantage or universal circuit optimality.

\section{Scope, limitations, and non-claims}

This paper does \emph{not}:
\begin{enumerate}[label=(\arabic*), leftmargin=1.5em]
\item optimize entangling gates;
\item define a full quantum compiler;
\item prove universal superiority over all baseline compilers;
\item claim hardware speedups;
\item claim overall circuit optimality; or
\item solve the general single-qubit synthesis problem in arbitrary gate sets.
\end{enumerate}

It defines and studies a local optimization pass for single-qubit segments.

\section{Conclusion}

The quaternionic model is operationally useful because it turns adjacent
single-qubit fusion into ordinary multiplication in a canonical representation.
Once paper 2's deterministic $\uoneq$ representation is in place, paper 3's
fusion rule follows naturally:
\[
\Fuse(\mathrm{seg})=\canon(q_n\cdots q_1).
\]
That rule preserves exact segment action modulo global phase, yields a unique
canonical fused form, and supports an honest reconstruction layer that can emit
identity, named gates, axis-aligned rotations, or generic primitives as needed.

On the shipped reference corpus, the local pass reduces 36 input single-qubit
gates to 11 reconstructed outputs while preserving exact equivalence under the
authoritative canonical-quaternion check. This is not a claim of universal
compiler optimality. It is a concrete result about one local method: a
geometry-native optimization pass for single-qubit segments.

\appendix

\section{Quaternion product formula}

For
\[
q_1=(w_1,x_1,y_1,z_1),\qquad q_2=(w_2,x_2,y_2,z_2),
\]
the product used by the fusion pass is
\[
q_2 q_1=
\bigl(
w_2 w_1 - x_2 x_1 - y_2 y_1 - z_2 z_1,\,
w_2 x_1 + x_2 w_1 + y_2 z_1 - z_2 y_1,
\]
\[
w_2 y_1 - x_2 z_1 + y_2 w_1 + z_2 x_1,\,
w_2 z_1 + x_2 y_1 - y_2 x_1 + z_2 w_1
\bigr).
\]

\section{Benchmark harness notes}

The reference harness included with this paper package uses only the Python
standard library. It translates each benchmark segment into quaternions, fuses
it by \eqref{eq:fuse}, applies reconstruction, and verifies equivalence by
re-translating the reconstructed output and comparing canonical quaternions.

\section{Companion figures and code}

Publication-oriented SVG figures are provided in
\texttt{artifacts/figures/}. The benchmark harness and measured results are
provided in \texttt{artifacts/code/benchmark\_gate\_fusion.py} and
\texttt{artifacts/code/benchmark\_results.json}.

\bibliographystyle{plainnat}
\bibliography{references}

\end{document}
